\question[3]
Gegeben ist ein Graph mit einem negativen Kantengewicht.

\includegraphics[height=2.5cm]{\pfad/dijkstra_08.png}
\begin{parts}
\part In welcher Reihenfolge werden die Knoten endgültig markiert, wenn der Algorithmus von Dijkstra bei a startet?
\part Welchen Wert hat dist(b) am Ende, welche Länge hat der kürzeste Weg von a nach b wirklich?
\part Warum liefert der Algorithmus hier ein falsches Ergebnis?
\end{parts}
\begin{solutionbox}{3.5cm}
a) a, b, c \\
b) dist(b) ist 2, der kürzeste Weg a-c-b hat die Länge 1. \\
c) b wird endgültig, weil er der billigste vorläufige Knoten ist. Das ist nur richtig, wenn kein Umweg billiger sein kann. Die negative Kante macht den späteren Umweg billiger.
\end{solutionbox}
